\clearpage
\fancyhead[L]{\small GAUSS AN BODE.}
\hypertarget{native-433-part-lower}{}
\begin{center}
{\Large GAUSS AN BODE.}

\rule{0.6\linewidth}{0.4pt}

Astronomisches Jahrbuch für 1823. Seite 160..161. Berlin 1820.
\end{center}

\hfill Göttingen den 21. März 1820.

\begin{center}
\emph{Saturn.}
\end{center}

\begin{center}
\begin{tabular}{r|c|c}
1819 Sept. 4 & $13^{\mathrm u}\,5^{\mathrm m}\,43^{\mathrm s}9$ & $359^\circ\,40'\,2''7$ \\
5 & $13\;\;1\;\;22.0$ & $36\;\;\;\;2.7$ \\
6 & $12\;57\;10.7$ & $32\;\;\;\;9.6$ \\
7 & $12\;52\;58.8$ & $28\;\;10.3$ \\
8 & $12\;48\;46.1$ & $24\;\;\;\;4.0$ \\
9 & $12\;44\;34.4$ & $20\;\;\;\;0.4$ \\
10 & $12\;40\;22.2$ & $15\;\;55.6$ \\
11 & $12\;36\;10.1$ & $11\;\;51.7$ \\
12 & $12\;31\;57.5$ & $\;\;7\;\;41.2$ \\
\hline
1819 Sept. 13 & $12^{\mathrm u}\,27\mathrm{m}\,45^{\mathrm s}0$ & $359^\circ\,3'\,30''4$ \\
14 & $12\;23\;32.3$ & $358\;59\;18.1$ \\
15 & $12\;19\;19.8$ & $55\;\;\;\;9.0$ \\
16 & $12\;15\;\;7.0$ & $50\;\;54.4$ \\
19 & $12\;\;2\;28.6$ & $38\;\;10.6$ \\
21 & $11\;54\;\;3.0$ & $29\;\;42.9$ \\
24 & $11\;41\;24.2$ & $15\;\;56.4$ \\
28 & $11\;24\;33.4$ & $\;\;0\;\;\;\;5.1$ \\
Oct. 1 & $11\;11\;55.9$ & $357\;47\;36.9$ \\
\end{tabular}
\end{center}

\Needspace{19\baselineskip}
\SourceLeaf{434}
\begin{center}
\textit{Vesta.}
\end{center}

\begin{center}
\begin{tabular}{r|c|c}
1819 Sept. 7 & $13^{\mathrm u}32^{\mathrm m}52^{\mathrm s}0$ & $9^\circ\,28'\,7''2$ \\
9 & $13\ 23\ 32.6$ & $9\ 6\ 10.5$ \\
10 & $13\ 18\ 51.1$ & $8\ 54\ 44.7$ \\
11 & $13\ 14\ 8.6$ & $8\ 43\ 1.8$ \\
12 & $13\ 9\ 25.1$ & $8\ 31\ 6.9$ \\
15 & $12\ 55\ 18.7$ & $7\ 53\ 49.8$ \\
\hline
1819 Sept. 16 & $12^{\mathrm u}50^{\mathrm m}21^{\mathrm s}4$ & $7^\circ\,40'\,58''5$ \\
19 & $12\ 35\ 55.5$ & $7\ 1\ 18.0$ \\
21 & $12\ 26\ 15.9$ & $6\ 34\ 15.9$ \\
24 & $12\ 11\ 43.4$ & $5\ 52\ 58.1$ \\
28 & $11\ 52\ 18.1$ & $4\ 57\ 24.3$ \\
Oct. 1 & $11\ 37\ 45.3$ & $4\ 16\ 1.2$ \\
\end{tabular}
\end{center}

\Needspace{12\baselineskip}
\begin{center}
\textit{Pallas.}
\end{center}

\begin{center}
\begin{tabular}{r|c|c}
1820 Jan. 2 & $12^{\mathrm u}0^{\mathrm m}56^{\mathrm s}9$ & $101^\circ\,44'\,50''4$ \\
9 & $11\ 27\ 31.7$ & $100\ 16\ 8.8$ \\
14 & $11\ 3\ 53.7$ & $99\ 16\ 26.5$ \\
\hline
1820 Jan. 15 & $10^{\mathrm u}59^{\mathrm m}12^{\mathrm s}7$ & $99^\circ\,5'\,5''1$ \\
22 & $10\ 26\ 58.4$ & $97\ 54\ 9.0$ \\
31 & $9\ 47\ 19.9$ & $96\ 50\ 9.0$ \\
\end{tabular}
\end{center}

\Needspace{12\baselineskip}
\begin{center}
\textit{Mars.}
\end{center}

\begin{center}
\begin{tabular}{r|c|c}
1820 Jan. 17 & $12^{\mathrm u}7^{\mathrm m}30^{\mathrm s}8$ & $118^\circ\,10'\,41''1$ \\
22 & $11\ 39\ 24.9$ & $116\ 3\ 44.8$ \\
23 & $11\ 33\ 50.8$ & $115\ 39\ 8.0$ \\
25 & $11\ 22\ 56.7$ & $114\ 50\ 55.9$ \\
\hline
1820 Febr. 8 & $10^{\mathrm u}9^{\mathrm m}43^{\mathrm s}1$ & $110^\circ\,19'\,58''3$ \\
9 & $10\ 4\ 52.7$ & $110\ 6\ 18.6$ \\
\end{tabular}
\end{center}

\par\indent Die Beobachtungen der \textit{Ceres} habe ich noch nicht reducirt. Es ist hiebei noch zu bemerken,
dass die Beobachtungen vom 17. Jan. an (incl.) mit \textsc{Bessels} neuer Tafel für die scheinbaren
geraden Aufsteigungen der \textsc{Maskelyn}ischen Sterne reducirt sind.

{\rightskip=0pt plus 3em
\par\indent Ueber den \textsc{Reichenbach}schen Meridiankreis, welcher seit dem 21. Februar in täglichem
Gebrauch ist, ein andermal.\par}
